% Note d'intention du CIR 2 — pendant de « Note_intention CIR1.pdf »
% (Fichierstravail/), dont elle reprend le plan et la mise en page.
% Compilation : latexmk -lualatex note-intention-CIR2.tex
\documentclass[11pt,a4paper]{article}
\usepackage[french]{babel}
\usepackage[a4paper,top=20mm,bottom=22mm,left=22mm,right=22mm,
            headheight=14pt,headsep=8mm,footskip=12mm]{geometry}
\usepackage{xcolor}
\usepackage{tikz}
\usepackage{tcolorbox}
\usepackage{enumitem}
\usepackage{booktabs}
\usepackage{tabularx}
\usepackage{titlesec}
\usepackage{fancyhdr}
\usepackage{lastpage}
\usepackage[hidelinks]{hyperref}

\definecolor{bleu}{HTML}{1F5FA6}
\definecolor{gristexte}{HTML}{555555}

\setlength{\parindent}{0pt}
\setlength{\parskip}{4pt}

\titleformat{\section}{\color{bleu}\Large\bfseries}{\thesection.}{0.4em}{}
  [{\color{bleu!45}\titlerule[0.4pt]}]
\titlespacing*{\section}{0pt}{10pt}{3pt}

\setlist[itemize]{label={\color{bleu}\small\textbullet},
  topsep=2pt,itemsep=2pt,parsep=0pt,leftmargin=1.6em}

\pagestyle{fancy}
\fancyhf{}
\fancyhead[L]{\color{bleu}\textbf{CIR2} — Certificat d'Initiation à la Radio, niveau 2}
\fancyhead[R]{\color{gristexte}Note d'intention}
\renewcommand{\headrule}{{\color{bleu}\hrule height 0.8pt}}
\fancyfoot[L]{\scriptsize\color{gristexte}F4JWI — F4JTL}
\fancyfoot[C]{\color{gristexte}Page \thepage{} / {\color{bleu}\pageref{LastPage}}}
\fancyfoot[R]{\scriptsize\color{gristexte}Licence CC BY 4.0 — Version d'octobre 2026}

% Pictogramme radio du bandeau, comme dans la note du CIR 1
\newcommand{\pictoradio}{%
  \begin{tikzpicture}[white, line width=1.6pt, line cap=round]
    \fill (0,0) circle (2.6pt);
    \foreach \r in {0.25,0.45,0.65}
      \draw (25:\r) arc (25:155:\r);
  \end{tikzpicture}}

\begin{document}

\begin{tcolorbox}[colback=bleu, colframe=bleu, arc=2mm,
                  left=4mm, right=4mm, top=2mm, bottom=2mm]
  \color{white}
  \begin{minipage}[c]{0.8\linewidth}
    {\LARGE\bfseries CIR2 — Note d'intention}\par\smallskip
    {\color{white!85!bleu}Présentation du projet \quad|\quad Collège — 4\textsuperscript{e} et 3\textsuperscript{e}}
  \end{minipage}\hfill
  \begin{minipage}[c]{0.15\linewidth}\raggedleft\pictoradio\end{minipage}
\end{tcolorbox}

\section{Le CIR2}

Le \textbf{CIR2} est le \textbf{second niveau du Certificat d'Initiation à la
Radio}. Il reprend, en les approfondissant, les notions découvertes au CIR1 et
conduit les élèves jusqu'à \textbf{la pratique} : écouter et identifier une
station lointaine, mesurer dans un circuit, construire et régler une antenne,
décoder un mode numérique, préparer puis tenir un vrai contact radio.

Le CIR2 s'adresse aux \textbf{élèves de 4\textsuperscript{e} et de
3\textsuperscript{e}}. Il ne suppose pas d'avoir suivi le CIR1 : les notions de
base sont reprises dès les premières séances. Il constitue la deuxième étape
d'un parcours, après le \textbf{CIR1} (6\textsuperscript{e} et
5\textsuperscript{e}) et avant la \textbf{préparation à la licence de
radioamateur} au lycée. Chaque séance est rattachée, dans le référentiel, au
programme du certificat d'opérateur des services d'amateur.

Le CIR2 est \textbf{adapté du cours en ligne 50ohm.de}, écrit par le
50ohm.de-Autorenteam et le référat AJW du DARC e.\,V., l'association allemande
des radioamateurs, et publié sous licence libre. Il n'en est pas une
traduction : le texte a été réécrit pour des collégiens français, le droit
français y est le texte courant, et les questions d'évaluation sont neuves.

Ce travail a été conçu par \textbf{Pierre LE GALL F4JWI} et \textbf{Jérôme
DIEUDONNE F4JTL}. Les auteurs sont \textbf{professeurs de mathématiques}. La
rédaction du CIR2 a été réalisée avec l'assistance d'une IA. La conception, les
choix pédagogiques et la relecture sont ceux des auteurs.

Le CIR2 est diffusé sous licence \textbf{Creative Commons Attribution 4.0
International (CC BY 4.0)}, comme le cours dont il est issu : il peut être
librement copié, modifié et redistribué, à condition de citer ses auteurs et
ceux du cours d'origine.

\section{Par des profs, pour des profs}

Le CIR2 a été pensé pour simplifier le travail des enseignants. En particulier :
\begin{itemize}
  \item il est prévu pour être traité \textbf{en une année scolaire}, à raison
    d'\textbf{une séance de 1\,h\,30 par semaine} ;
  \item chaque séance suit la même trame : \textbf{25 minutes de cours},
    \textbf{50 minutes de manipulation} par groupes de deux ou trois, puis
    \textbf{15 minutes de bilan} ;
  \item un \textbf{radioamateur} est présent à chaque séance : il fait les
    démonstrations à l'antenne et apporte son expérience ; l'enseignant n'a
    pas besoin d'être lui-même radioamateur ;
  \item les séances suivent une progression, mais seuls quelques blocs sont
    liés (les antennes, séances 14 à 16 ; le contact, séances 21 et 22) : les
    autres peuvent être déplacées ;
  \item le \textbf{matériel} est décrit dans une mallette type, en trois
    niveaux d'équipement ; les exercices de trafic entre élèves se font avec
    des \textbf{postes PMR446}, utilisables sans licence ;
  \item l'évaluation finale se fait \textbf{en ligne} et se corrige seule ;
  \item à la radio, les élèves ne donnent que leur \textbf{prénom} et leur
    \textbf{ville} : jamais leur nom de famille ni leur adresse ; les exercices
    se font sous des indicatifs fictifs.
\end{itemize}

\section{Le programme}

Le cours du CIR2 compte \textbf{24 séances de 1\,h\,30}, soit 36 heures,
réparties en \textbf{cinq domaines} : \textbf{A}~découvrir la radio,
\textbf{B}~électricité et électronique, \textbf{C}~ondes, spectre et
propagation, \textbf{D}~modulations, antennes et station, \textbf{E}~trafic,
règles et sécurité.

\newpage

{\small
\begin{tabularx}{\linewidth}{@{}c c >{\raggedright\arraybackslash}p{0.33\linewidth} >{\raggedright\arraybackslash}X@{}}
\toprule
\textbf{Séance} & \textbf{Dom.} & \textbf{Titre} & \textbf{Manipulation} \\
\midrule
1  & A & La radio autour de nous & première écoute, récepteur et WebSDR \\
2  & E & Se présenter à l'antenne & épeler et noter des indicatifs en PMR \\
3  & E & Un contact radio & contact simulé en PMR, carnet de trafic \\
4  & B & Tension et courant & circuit pile-lampe, mesures au multimètre \\
5  & B & La résistance et la loi d'Ohm & vérifier la loi d'Ohm sur trois résistances \\
6  & B & Puissance et énergie & mesurer une puissance, calculer une autonomie \\
7  & C & Du courant alternatif à l'onde & visualiser une sinusoïde \\
8  & C & Le spectre et les bandes & dresser sa carte du spectre à la clé SDR \\
9  & C & La propagation à vue & mesurer la portée d'un poste PMR \\
10 & C & Les ondes courtes et l'ionosphère & écouter balises et stations lointaines \\
11 & D & Moduler : transporter la voix & reconnaître AM, FM et SSB \\
12 & D & Le Morse & envoyer et recevoir son prénom \\
13 & D & Les modes numériques et l'image & décoder du FT8 et une image SSTV \\
14 & D & Les antennes : principes & calculer et couper un dipôle 2\,m \\
15 & D & Construire son antenne & en faire une antenne directive, comparer \\
16 & D & Câbles, connecteurs et réglage & régler l'antenne à l'analyseur \\
17 & B & Le récepteur & prise en main d'un poste en réception \\
18 & B & L'émetteur & démonstration sur charge fictive \\
19 & B & L'atelier de soudure & souder un petit kit \\
20 & A & La radio en action & chasse au renard dans l'établissement \\
21 & E & Préparer son contact & répétition générale en PMR, script en main \\
22 & E & Au micro & un vrai contact, aux côtés du radioamateur \\
23 & E & Les règles, et après ? & examen blanc corrigé \\
\midrule
24 & — & L'évaluation du CIR2 & évaluation, correction, remise des attestations \\
\bottomrule
\end{tabularx}
}

\section{Les documents fournis}

Pour l'ensemble du parcours, sont fournis :
\begin{itemize}
  \item un \textbf{manuel de l'élève}, un chapitre par séance, avec ses
    questions d'entraînement ;
  \item un \textbf{livre du formateur}, une fiche par séance : ce que l'élève
    doit retenir, préparation, déroulé minuté, conduite de la manipulation,
    difficultés fréquentes, pour aller plus loin, corrigé ;
  \item des \textbf{fiches d'activité}, une par séance, en deux versions :
    l'une pour l'élève, l'autre pour l'encadrant ;
  \item une \textbf{banque de 230 questions}, dix par séance, exportable vers
    une plateforme d'enseignement ;
  \item un \textbf{référentiel} : connaissances et savoir-faire attendus, et
    leur rattachement au programme du certificat d'opérateur ;
  \item un \textbf{livret de l'élève}, où sont validés les savoir-faire, et
    une \textbf{mallette type}, chiffrée.
\end{itemize}

\section{L'évaluation du CIR2}

L'évaluation se passe \textbf{en ligne}, sur une plateforme de type Moodle ou
éléa : \textbf{40 questions à choix multiple} en \textbf{une heure}, soit
\textbf{8 questions par domaine}, tirées au hasard dans la banque du CIR2. Un
examen blanc corrigé est prévu à la séance 23.

Le CIR2 est obtenu avec \textbf{au moins 20 bonnes réponses sur 40}, dont
\textbf{au moins 3 dans chaque domaine}. Les savoir-faire pratiques sont
validés au fil des séances dans le livret de l'élève, sans note.

\section{Remerciements}

Les auteurs remercient le \textbf{Réseau des Émetteurs Français (REF)} pour
son soutien, Freddy \textbf{ON6FS} pour sa participation aux premières étapes
du projet, ainsi que \textbf{F4LDR} et \textbf{F4MHI} pour leurs relectures.

\end{document}
